\documentclass[10pt,a4paper]{amsart}
\usepackage[margin=19mm]{geometry}
\usepackage{amsmath,amssymb,mathtools}
\usepackage[scr=rsfs,cal=euler]{mathalfa}
\usepackage{xfrac}
\usepackage[unicode,hidelinks,bookmarks=false]{hyperref}
\newcommand{\ie}{i.e.\ }
\newcommand{\cf}{cf.\ }
\newcommand{\resp}{respectively\ }
\newcommand{\Subsection}{\subsection}
\providecommand{\nobreakdash}{\nobreak\hbox{-}}
\setlength{\emergencystretch}{2em}
\multlinegap0pt
\usepackage{amssymb}
\usepackage{dsfont}
\usepackage{mathtools}
\usepackage{csquotes}
\usepackage{bm}
\usepackage[shortlabels]{enumitem}
\usepackage{xcolor}
\usepackage{cancel}
\usepackage{float}
\usepackage{tikz}
\newtheorem{theorem}{Theorem}
\newtheorem{proposition}{Proposition}
\title{Polyconvexity for Cosserat nonlinear elasticity and nonlinear couple--stress theory}
\author{Diana Ciotir, Ionel-Dumitrel Ghiba, Patrizio Neff}
\date{Source: arXiv:2608.23072v1 (CC BY 4.0)}
\begin{document}
\maketitle

\noindent\emph{Source and licence.} Polyconvexity for Cosserat nonlinear elasticity and nonlinear couple--stress theory. Version arXiv:2608.23072v1 (CC BY 4.0), distributed by arXiv under the Creative Commons Attribution 4.0 International licence, http://creativecommons.org/licenses/by/4.0/, as recorded in the arXiv licence metadata for this exact version and shown on the abstract page https://arxiv.org/abs/2608.23072v1. Full licence notice: \texttt{licenses/arxiv-2608.23072-CC-BY-4.0.txt}.\par\smallskip

\noindent\textit{Editorial scope.} Counted unit: the article's theorem thm:elastic3 (source0.tex lines 2445--2448). The article's complete original proof of it is delivered with it (source0.tex lines 2450--2484, one \texttt{proof} environment of 35 lines). No mathematics is written, repaired or re-derived: the statement and the proof are the article's own lines, in the article's own notation, and no retained line is substituted or re-flowed.  Retained uncounted as the source-local context the proof needs: lines 2159--2167; lines 2376--2418.  Nothing was omitted from any retained span, so the package carries no omission record.  No retained line contains a citation, so the wrapper carries no bibliography and no bibliography entry.  No retained line contains a figure environment, an image-inclusion command or drawing code, so no image is delivered and the package declares no asset dependency.  Editorial formatting only, in addition: an AMS \texttt{amsart} preamble that restates the article's own declarations and macros used by the retained lines, namely the environments \texttt{theorem}, \texttt{proposition} on separate counters, together with the standard packages those lines need.  The retained environments print in this document as Theorem 1, Proposition 1 in order of appearance, whereas the article prints the same items with the numbering of its own full document.  The complete native source of the article as distributed in the arXiv e-print for this exact version is bundled unchanged as \texttt{sources/208390/source0.tex}.
\begin{theorem}[Existence for the auxiliary lifted problem]
\label{thm:elastic2}
Assume that $\mathcal A_{\rm Cosserat}$ contains a pair of finite energy. If either H1)--H4) or H1')--H4') holds, then
$
\inf_{(\varphi, \overline R )\in\mathcal A_{\rm Cosserat}}
\widehat{\mathcal I}(\varphi, \overline R )
$
admits a minimizer.
\end{theorem}

\begin{proposition}[Polar recovery and equivalence of admissible classes]
\label{prop:polar-equivalence}
Define
\begin{align}
\mathcal A_{\rm pol}
:=
\Bigl\{
\varphi\in {\rm W}^{1,p}(\Omega;\mathbb R^3)\,|\,
&\operatorname{Tr}_{\Gamma_d}\varphi=\varphi^*,
\quad
\det{\rm D}\varphi>0\ \text{a.e.},
\notag\\[-1mm]
& R:=\operatorname{polar}({\rm D}\varphi)
\in  {\rm W}^{1,q}(\Omega;{\rm SO}(3)) 
\Bigr\}.
\end{align}
Then
$
\varphi\mapsto(\varphi,R)
$
is a bijection from $\mathcal A_{\rm pol}$ onto $\mathcal A_{\rm relax}$, and, therefore, the admissible classes $\mathcal A_{\rm pol}$ onto $\mathcal A_{\rm relax}$ are equivalent. For every corresponding pair,
$
R^T{\rm D}\varphi
=
\sqrt{({\rm D}\varphi)^T{\rm D}\varphi}
\in\operatorname{Sym}^{++}(3)
\text{a.e.}
$
Moreover, if
\begin{align}
\mathcal J(\varphi)
:=
\int_\Omega
W\left(x,
\sqrt{({\rm D}\varphi)^T{\rm D}\varphi},
R^T\operatorname{Curl}\overline R
\right)\,\mathrm dx,
\end{align}
then
$
\mathcal J(\varphi)=\mathcal I(\varphi,R).
$
\end{proposition}

\begin{theorem}[Existence in the constrained and polar formulations]
\label{thm:elastic3}
Assume that $\mathcal A_{\rm relax}$ contains a pair of finite energy. Under either H1)--H4) or H1')--H4'), the functional $\mathcal I$ admits a minimizer in $\mathcal A_{\rm relax}$. Equivalently, $\mathcal J$ admits a minimizer in $\mathcal A_{\rm pol}$.
\end{theorem}

\begin{proof}
Let $(\varphi_k, \overline R _k)\subset\mathcal A_{\rm relax}$ be a minimizing sequence. All compactness, identification of minors and lower-semicontinuity arguments from Theorem~\ref{thm:elastic2} apply without change. Thus, along a subsequence,
\begin{align}
\varphi_k&\rightharpoonup\varphi
&&\text{in }{\rm W}^{1,p}(\Omega;\mathbb R^3),
\notag\\
 \overline R _k&\to \overline  R 
&&\text{a.e. and strongly in every finite }{\rm L}^a,
\notag\\
\overline U_k:= \overline R _k^T{\rm D}\varphi_k
&\rightharpoonup
\overline U:= \overline R ^T{\rm D}\varphi
&&\text{in }{\rm L}^p(\Omega;\mathbb R^{3\times3}),
\end{align}
and the limit has positive determinant and finite energy.

For every $k$,
$
\overline U_k(x)\in\operatorname{Sym}^{+}(3)
$
almost everywhere. The set
\begin{align}
\left\{
V\in {\rm L}^p(\Omega;\mathbb R^{3\times3})\, |\,
V(x)\in\operatorname{Sym}^{+}(3)\ \text{a.e.}
\right\}
\end{align}
is convex and strongly closed, hence weakly closed. Therefore
$
\overline U= \overline R ^T{\rm D}\varphi\in\operatorname{Sym}^{+}(3)
$
almost everywhere. The trace condition also passes to the limit, so $(\varphi, \overline R )\in\mathcal A_{\rm relax}$. Weak lower semicontinuity gives minimality.

Proposition~\ref{prop:polar-equivalence} identifies $ \overline R $ with $\operatorname{polar}({\rm D}\varphi)$ and transfers the minimizer bijectively to $\mathcal A_{\rm pol}$.
\end{proof}

\end{document}
